%% Resumo
\begin{resumo}

O aumento do número de pessoas que vivem sozinhas, especialmente idosos e indivíduos com necessidades especiais, tem ampliado a preocupação com situações de emergência ocorridas sem assistência imediata. Nesse contexto, este trabalho apresenta o desenvolvimento do Linha Vital, um sistema mobile voltado ao monitoramento de inatividade e ao acionamento automatizado de contatos de confiança em situações de risco. O objetivo do projeto consiste em oferecer uma solução tecnológica acessível, não invasiva e capaz de contribuir para a segurança e autonomia de pessoas em condição de vulnerabilidade. A metodologia adotada envolveu pesquisa bibliográfica sobre monitoramento remoto de saúde, levantamento de requisitos, modelagem do sistema, definição de regras de negócio, elaboração de histórias de usuário e desenvolvimento da arquitetura da aplicação. A solução foi estruturada utilizando a linguagem Kotlin para o desenvolvimento mobile e backend, o framework Spring Boot para implementação dos serviços, o PostgreSQL para persistência de dados e a Oracle Cloud Infrastructure para hospedagem dos componentes da aplicação. O sistema realiza monitoramento por meio da identificação de períodos prolongados de inatividade, envio de notificações periódicas de verificação de bem-estar e acionamento automático de contatos previamente cadastrados quando determinadas condições de risco são identificadas. Além disso, disponibiliza recursos de acionamento manual de emergência e mecanismos voltados à acessibilidade. Os resultados obtidos demonstram a viabilidade técnica da proposta, evidenciando que a utilização de tecnologias amplamente consolidadas possibilita a construção de uma solução confiável, escalável e alinhada às necessidades do público-alvo. Conclui-se que o Linha Vital representa uma alternativa promissora para ampliar a segurança de indivíduos que vivem sozinhos, contribuindo para a redução do tempo de resposta em situações emergenciais e promovendo maior qualidade de vida por meio da integração entre tecnologia, acessibilidade e cuidado.

\vspace{\onelineskip}
\noindent
\textbf{Palavras-chave}: monitoramento remoto; tecnologia assistiva; acessibilidade; aplicativo móvel; saúde digital.
\end{resumo}